\section{Tests Beyond the Benchmark}
\label{sec:beyond}
This section tests how far the conclusions of Section~\ref{sec:results} carry over to conditions that the benchmark does not cover: larger budgets, feasible initial layouts, a real offshore site and an external benchmark with published layouts. The aim is to locate the limits of PSO-VNS, not to add further wins.

\subsection{Budget Scaling and Feasibility-Preserving Initialization}
\label{sec:feasinit}
\label{sec:budget}
%% generated by mpce_results.py -- do not edit by hand
\begin{table}[!t]
\centering
\caption{Six Largest Benchmark Cases: Mean Wake Loss (\%) and Feasible Runs (\%) at 6,030 Calls with Random and Feasibility-Preserving Initialization, and Average Rank at 6,030, 30,030 and 120,030 Calls (Random Initialization; 30 Seeds)}
\label{tab:feasbudget}
\scriptsize\setlength{\tabcolsep}{2.4pt}
\begin{tabular}{lccccccc}
\toprule
& \multicolumn{2}{c}{Random init.} & \multicolumn{2}{c}{Feasible init.} & \multicolumn{3}{c}{Avg.\ rank} \\
\cmidrule(lr){2-3}\cmidrule(lr){4-5}\cmidrule(lr){6-8}
Method & Loss & Feas. & Loss & Feas. & 6,030 & 30,030 & 120,030 \\
\midrule
PSO-VNS & 4.084 & 100.0 & 3.938 & 100.0 & 1.50 & 1.17 & 2.17 \\
SSA-VNS & 4.678 & 95.0 & \TBD{} & \TBD{} & 4.33 & 4.17 & 4.17 \\
LX-SSA-VNS & 4.802 & 98.9 & \TBD{} & \TBD{} & 5.17 & 5.67 & 5.83 \\
RS-VNS & 4.830 & 93.3 & \TBD{} & \TBD{} & 4.67 & 6.33 & 4.67 \\
VNS & 4.950 & 100.0 & \TBD{} & \TBD{} & 6.50 & 4.83 & 3.17 \\
PSO & 4.181 & 100.0 & \TBD{} & \TBD{} & 1.83 & 2.50 & 5.83 \\
SSA & 4.828 & 80.0 & \TBD{} & \TBD{} & 8.33 & 8.67 & 8.83 \\
LX-SSA & 5.379 & 77.8 & \TBD{} & \TBD{} & 8.67 & 9.83 & 10.00 \\
DE & -- & 1.1 & \TBD{} & \TBD{} & 10.00 & 5.17 & 3.33 \\
MS-SLSQP & 4.522 & 99.4 & \TBD{} & \TBD{} & 4.00 & 6.67 & 7.00 \\
\bottomrule
\multicolumn{8}{l}{\TBD{pending: feas}}
\end{tabular}
\end{table}

\begin{figure*}[!t]
\centering
\pipefig{0.7\textwidth}{figures_mpce/budget_scaling.pdf}{7.4cm}
\caption{Budget scaling: mean wake loss of the six largest benchmark cases and mean AEP loss of the Horns Rev~1 block (feasible runs; methods with at least half of their runs feasible) versus the budget (6,030, 30,030 and 120,030 calls, random initialization; 10 seeds for Horns Rev~1 at 120,030 calls).}
\label{fig:budget}
\end{figure*}

\emph{Budget.} With 5 and 20 times the budget, \NBudgetPhrase{} on the six largest benchmark cases (\NBudRankPSOVNSSixK, \NBudRankPSOVNSThirtyK{} and \NBudRankPSOVNSOneTwentyK; Table~\ref{tab:feasbudget}, Fig.~\ref{fig:budget}), and its mean wake loss falls from \NBudLossPSOVNSSixK\% to \NBudLossPSOVNSThirtyK\% and \NBudLossPSOVNSOneTwentyK\%. The mean wake loss of PSO minus that of PSO-VNS is \NBudGainPSOSixK, \NBudGainPSOThirtyK{} and \NBudGainPSOOneTwentyK{} percentage points at the three budgets (per-case values in Table~\ref{S-tab:budget-cases}). The order of the other methods, however, depends on the budget: the average rank of PSO drops from \NBudRankPSOSixK{} to \NBudRankPSOOneTwentyK, whereas VNS (\NBudRankVNSSixK{} to \NBudRankVNSOneTwentyK) and DE move up, and at 120,030 evaluations the mean wake losses of PSO-VNS, VNS and SSA-VNS (\NBudLossPSOVNSOneTwentyK\%, \NBudLossVNSOneTwentyK\% and \NBudLossSSAVNSOneTwentyK\%) are close. A comparison at a single budget can thus order the baselines differently; in our tests, only the first place of PSO-VNS is stable.
% CHECK-FINAL [C26]: feasbudget.rank.6030: PSOBV best (column 6,030; the 30,030 / 120,030 statements are generated by \NBudgetPhrase)
% verified 2026-09-28 from mpce_summary.json (feasbudget.rank, feasbudget.mean_loss): PSOC rank 1.83 -> 5.83, BVNS 6.50 -> 3.17, DE 10.00 -> 3.33 (6,030 -> 120,030); no CHECK-FINAL condition covers this sentence

\emph{Feasible starts.} With feasibility-preserving initialization, \NFeasInitPhrase. The best method is then VNS alone: \NFbFeasBestLoss{} has the lowest mean wake loss and \NFbFeasBestRank{} the best average rank, whereas PSO-VNS ranks \NFbFeasRankPSOVNS{} on average (mean wake loss \NFbRandLossPSOVNS\% with random and \NFbFeasLossPSOVNS\% with feasible starts; per-case values in Table~\ref{S-tab:feasinit-cases}); RS-VNS was not run with this initialization (Section~\ref{sec:setup-beyond}). PSO and DE never improve on the best initial layout in these runs: since the turbine labels are arbitrary, velocity and difference moves that combine two feasible layouts almost always make turbines overlap, the feasibility-first comparison rejects them, and the global-best particle of PSO does not move. The single-coordinate moves of VNS are not affected, so a feasible start favors the VNS-based methods, whereas the penalty-guided search from random starts favors PSO. The advantage of PSO-VNS therefore depends on the initialization.
% verified 2026-09-28 from mpce_feas_s0of1.csv: the convergence curves of PSO and DE show no improvement after the initial population in any of the 210 feasible-initialization runs (identical final layouts for both).
% verified 2026-09-28 from mpce_summary.json (feasbudget.rank_feasible_init): BVNS 1.00 (first), PSOBV 2.83; no CHECK-FINAL condition covers "VNS alone best"

\subsection{Horns Rev~1 16-Turbine Block}
\label{sec:hornsrev-site}
%% generated by mpce_results.py -- do not edit by hand
\begin{table}[!t]
\centering
\caption{Horns Rev~1 16-Turbine Block (Wake-Free AEP 149.26 GWh/yr): Mean AEP (GWh/yr) over the Feasible Runs of 30 Seeds at 6,030 Calls, Feasible Runs, Holm-Adjusted Wilcoxon $p$ of PSO-VNS vs.\ Each Method, and Mean AEP Loss (\%) at 6,030 Calls with Random (R) and Feasibility-Preserving (F) Initialization and at 30,030 and 120,030 Calls (R; 10 Seeds at 120,030); Superscript: Feasible Runs When Not All Runs Are Feasible}
\label{tab:hr-site}
\scriptsize\setlength{\tabcolsep}{2.2pt}
\begin{tabular}{lccccccc}
\toprule
& \multicolumn{3}{c}{6,030 calls, R} & \multicolumn{4}{c}{Loss (\%)} \\
\cmidrule(lr){2-4}\cmidrule(lr){5-8}
Layout / method & AEP & Feas. & $p_{\rm Holm}$ & 6k R & 6k F & 30k & 120k \\
\midrule
Installed layout & 139.51 & -- & -- & 6.53 & -- & -- & -- \\
\midrule
PSO-VNS & 138.73 & 30/30 & -- & 7.05 & 6.96 & 6.46 & \TBD{} \\
SSA-VNS & 137.98 & 30/30 & $6.7\times10^{-4}$ & 7.56 & \TBD{} & \TBD{} & \TBD{} \\
LX-SSA-VNS & 137.77 & 30/30 & $6.3\times10^{-4}$ & 7.70 & \TBD{} & \TBD{} & \TBD{} \\
RS-VNS & 137.96 & 28/30 & $2.2\times10^{-4}$ & 7.57$^{28}$ & \TBD{} & \TBD{} & \TBD{} \\
VNS & 137.67 & 30/30 & $2.2\times10^{-4}$ & 7.77 & \TBD{} & \TBD{} & \TBD{} \\
PSO & 138.06 & 19/30 & $1.4\times10^{-5}$ & 7.51$^{19}$ & \TBD{} & \TBD{} & \TBD{} \\
SSA & 136.66 & 30/30 & $2.8\times10^{-7}$ & 8.44 & \TBD{} & \TBD{} & \TBD{} \\
LX-SSA & 136.46 & 28/30 & $1.5\times10^{-7}$ & 8.58$^{28}$ & \TBD{} & \TBD{} & \TBD{} \\
DE & -- & 0/30 & $1.7\times10^{-8}$ & -- & \TBD{} & \TBD{} & \TBD{} \\
MS-SLSQP & 136.80 & 30/30 & $1.5\times10^{-7}$ & 8.35 & \TBD{} & \TBD{} & \TBD{} \\
\bottomrule
\end{tabular}
\end{table}


Table~\ref{tab:hr-site} gives the results for the 16-turbine block. Under our model, the installed layout, a regular $7D$ grid, reaches \NHRInstalled~GWh/yr (wake loss \NHRInstalledLoss\%). At 6,030 evaluations, PSO-VNS finds a feasible layout in all 30 runs and has the highest mean AEP, \NHRMeanPSOVNS~GWh/yr (SD \NHRSDPSOVNS), significantly higher than that of every other method ($p_{\rm Holm}\le\NHRMaxP$); this mean is still below the installed layout, which is exceeded by its best run (\NHRBestPSOVNS~GWh/yr) and \NHRAbovePSOVNS{} of its 30 runs; \NHRAboveOthersPhrase. At 30,030 evaluations, the mean AEP of PSO-VNS, \NHRMeanThirtyKPSOVNS~GWh/yr (loss \NHRLossThirtyKPSOVNS\%), exceeds that of the installed layout, as do \NHRAboveThirtyKPSOVNS{} of the 30 runs; at 120,030 evaluations the mean loss is \NHRLossOneTwentyKPSOVNS\%. At this largest budget, however, VNS, RS-VNS and DE reach lower mean losses than PSO-VNS, so the site confirms the budget dependence seen above. From random starts, PSO finds a feasible layout inside the parallelogram in only \NHRFeasPSO{} runs; with feasibility-preserving initialization it does so in \NHRFeasFeasInitPSO{} runs, and at 30,030 and 120,030 evaluations in \NHRFeasThirtyKPSO{} and \NHRFeasOneTwentyKPSO{} runs. The best layouts are shown in Fig.~\ref{fig:layouts}. The comparison with the installed layout uses our wake model and boundary and does not assess the design of the real farm.
% CHECK-FINAL [C19]: hr16: PSOBV highest mean AEP, all p_holm < 0.05, best > installed
% CHECK-FINAL [C20]: hr16.methods.PSOC.feasible < runs ("only")
% CHECK-FINAL [C25]: hr16.loss_by_setting.PSOBV: 30030R and 120030R loss < 6030R ("falls")
% verified 2026-09-28 from mpce_numbers.tex / mpce_tab_hr16.tex: mean 6,030 (138.73) < installed (139.51) < mean 30,030 (139.62); at 120,030 BVNS 6.12, RSVNS 6.24, DE 6.34 < PSOBV 6.45; no CHECK-FINAL condition covers these two statements

\subsection{IEA37 Case Study~1}
\label{sec:iea37}
%% generated by mpce_results.py -- do not edit by hand
\begin{table}[!t]
\centering
\caption{IEA37 Case Study~1, 16- and 36-Turbine Scenarios: AEP in MWh of the Best Run of Each Method at 6,030 and 30,030 Calls (30 Seeds), Compared with the Example Layout and the Best Feasible Published Layout~\cite{Baker2019,IEA37repo}}
\label{tab:iea37}
\scriptsize\setlength{\tabcolsep}{2.5pt}
\begin{tabular}{lcccc}
\toprule
& \multicolumn{2}{c}{16 turbines ($r=1300$~m)} & \multicolumn{2}{c}{36 turbines ($r=2000$~m)} \\
\cmidrule(lr){2-3}\cmidrule(lr){4-5}
Method / source & 6,030 & 30,030 & 6,030 & 30,030 \\
\midrule
Example layout & \multicolumn{2}{c}{366{,}941.6} & \multicolumn{2}{c}{737{,}883.1} \\
Best feasible published$^{a}$ & \multicolumn{2}{c}{418{,}924.4} & \multicolumn{2}{c}{863{,}676.3} \\
\midrule
PSO-VNS & \TBD{} & \TBD{} & \TBD{} & \TBD{} \\
PSO & 403{,}753.1 & 412{,}702.4 & 787{,}988.7 & 830{,}256.4 \\
SSA-VNS & 403{,}251.3 & 406{,}924.6 & 766{,}803.3 & 819{,}930.5 \\
SSA & 395{,}150.2 & 401{,}459.3 & 750{,}707.4 & 802{,}538.2 \\
LX-SSA & 393{,}284.6 & 398{,}770.0 & 746{,}301.3 & 789{,}311.3 \\
DE & 357{,}391.3 & 397{,}402.7 & -- & -- \\
VNS & 400{,}169.3 & 404{,}076.0 & 785{,}791.3 & 824{,}230.2 \\
MS-SLSQP & 405{,}433.8 & 410{,}349.5 & 748{,}571.1 & 833{,}365.0 \\
\bottomrule
\multicolumn{5}{p{0.92\columnwidth}}{$^{a}$Participant~4 in both scenarios (SNOPT with wake expansion continuation \TBD{verify against Baker et al.\ (2019)}). A higher submitted AEP places turbines up to 3.5~m outside the boundary and is not counted.}
\end{tabular}
\end{table}

\begin{figure*}[!t]
\centering
\pipefig{0.9\textwidth}{figures_mpce/layouts_iea37_hr16.pdf}{5.4cm}
\caption{Best layouts. Left and middle: IEA37 Case Study~1, 16 and 36 turbines, best PSO-VNS layout at the largest budget and best feasible published layout. Right: Horns Rev~1 16-turbine block, installed layout and best PSO-VNS layout at 6,030 calls.}
\label{fig:layouts}
\end{figure*}

Table~\ref{tab:iea37} and Fig.~\ref{fig:layouts} compare our methods with the layouts submitted to IEA37 Case Study~1, in which the participants chose their own optimizers and budgets~\cite{Baker2019}. None of our methods reaches the best feasible published layouts, which come from a gradient-based method run with a budget of its own choice (footnote of Table~\ref{tab:iea37}). At 30,030 evaluations, the relative AEP difference of the best PSO-VNS layout to the best feasible published layout is \NIEAGapSixteenThirtyK\% for 16 turbines and \NIEAGapThirtySixThirtyK\% for 36 turbines (mean over the 30 runs: \NIEAMeanGapSixteenThirtyK\% and \NIEAMeanGapThirtySixThirtyK\%; best run at 6,030 evaluations: \NIEAGapSixteenSixK\% and \NIEAGapThirtySixSixK\%). The best PSO-VNS layout would rank \NIEARankSixteenThirtyK{} among the \NIEAPubFeasSixteen{} feasible submitted 16-turbine layouts and \NIEARankThirtySixThirtyK{} among the \NIEAPubFeasThirtySix{} feasible 36-turbine layouts. The gap is larger for 36 turbines, and there the best of our layouts at 30,030 evaluations comes from \NIEABestOfThirtySixThirtyK, which also uses gradients. On this external benchmark, the metaheuristics compared here are therefore not competitive with the best gradient-based results; since the participants' budgets are not reported in a comparable form, the comparison is indicative only.
% verified 2026-09-28 from mpce_numbers.tex: gap 16/30k -1.33 % vs 36/30k -4.20 % ("larger for 36 turbines"); NIEABestOfThirtySixThirtyK = MS-SLSQP ("also uses gradients"); no CHECK-FINAL condition covers these statements

\subsection{Published Kusiak--Song Results and Computational Cost}
\label{sec:literature}
\label{sec:runtime}
A comparison with published results on the Kusiak--Song benchmark~\cite{Kusiak2010,Eroglu2012,Eroglu2013,Bansal2017} is given in Section~\ref{S-sec:S-literature} of the supplementary material \TBD{pending author verification of the PDFs}. At equal numbers of evaluations, all methods need a similar time (\NMsPerCallMin--\NMsPerCallMax~ms per evaluation on the machines that produced the respective runs; Table~\ref{S-tab:cost}); one 6,030-evaluation run of PSO-VNS takes \NSecPerRunMin--\NSecPerRunMax~s, depending on $N$.
